% ============================================================================
% Fichier  : partie4/P04_C14_web.tex
% Titre    : Tester une Application Web
% Auteur   : MINKA MI NGUIDJOÏ Thierry Emmanuel
% Version  : 1.0 -- novembre 2026
% ============================================================================
% Sources : OWASP Web Security Testing Guide (WSTG) ; OWASP Top 10:2021 ;
% RFC 9110 (sémantique HTTP) ; RFC 6265 (cookies). Portail de suivi de
% LiZbiZ et Opération LOGONE : cas fictifs. Deux séances (S17 et S18).
% ============================================================================

\chapter{Tester une Application Web}
\label{chap:web}
\chaptermark{Tester une application web}

\epigraph{%
    « Une application web fait confiance à ce que le navigateur lui envoie~;\\
    le testeur, lui, n'utilise pas forcément un navigateur. »%
}{Principe du cours}

\niveauChapitre{Intermédiaire}{Chapitre~\ref{chap:methodes}~; cours SEC3266,
chapitre sur la sécurité applicative.}

\begin{boxobjectifs}
\begin{itemize}
    \item Décrire une application web du point de vue du testeur~: requêtes,
          réponses, état, session.
    \item Utiliser le guide de test de l'OWASP (WSTG) pour structurer une
          campagne.
    \item Cartographier une application~: pages, paramètres, rôles.
    \item Tester l'authentification et la gestion de session.
    \item Lire un historique HTTP pour y repérer une faiblesse.
    \item Classer les constats selon le Top~10 de l'OWASP.
\end{itemize}
\end{boxobjectifs}

\bigskip

% ============================================================================
\section*{Question directrice : par où commencer le test d'un portail que l'on ne connaît pas~?}
% ============================================================================

Le portail de suivi des colis de LiZbiZ permet aux clients de voir leurs
expéditions et à 40~agents de les mettre à jour. M.~TAMO reçoit deux comptes
de test, un client et un agent, et l'adresse du portail de pré-production.
Il dispose de 48~heures. Ce chapitre montre comment cartographier, puis
tester méthodiquement, plutôt que d'essayer au hasard.

\medskip
\textbf{Plan du chapitre.}
\begin{enumerate}
    \item L'application vue par le testeur (section~\ref{sec:c14-modele}).
    \item Le guide de test de l'OWASP (section~\ref{sec:c14-wstg}).
    \item Cartographier (section~\ref{sec:c14-carto}).
    \item Authentification et session (section~\ref{sec:c14-auth}).
    \item Lire un historique HTTP (section~\ref{sec:c14-historique}).
    \item Classer selon le Top~10 (section~\ref{sec:c14-top10}).
\end{enumerate}

\separateur

% ============================================================================
\section{L'Application vue par le Testeur}
\label{sec:c14-modele}
% ============================================================================

\begin{boxdef}
\termecle{Requête HTTP}~: message envoyé par le client au serveur (méthode,
chemin, en-têtes, corps). \termecle{Réponse HTTP}~: message du serveur
(code, en-têtes, corps). \termecle{Session}~: état conservé entre requêtes,
le plus souvent par un cookie (RFC~6265).
\end{boxdef}

Le testeur observe et modifie chaque requête avant son envoi. Un serveur
qui se fie à un contrôle fait côté navigateur (champ masqué, menu grisé,
vérification en JavaScript) accepte une requête forgée qui ignore ce
contrôle. La règle fondatrice du test web~: \emph{tout ce qui vient du
client est contrôlable par le client}.

% ============================================================================
\section{Le Guide de Test de l'OWASP}
\label{sec:c14-wstg}
% ============================================================================

Le \termecle{Web Security Testing Guide} (WSTG) de l'OWASP organise le test
d'une application web en catégories~\cite{owaspwstg}. On ne les parcourt pas
toutes à chaque mission~; on choisit selon le périmètre et le temps.

\begin{tableaucours}{L{4.4cm}L{7.8cm}}{%
    \tableauHeader{Catégorie WSTG} &
    \tableauHeader{Ce qu'on y vérifie}
}
Collecte d'informations & Technologies, points d'entrée, pages cachées \\
Configuration & En-têtes, méthodes autorisées, messages d'erreur \\
Gestion d'identité et authentification & Création de compte, connexion,
réinitialisation \\
Gestion de session & Cookies, expiration, déconnexion \\
Contrôle d'accès & Un rôle peut-il atteindre les données d'un autre~? \\
Validation des entrées & Injections, scripts, chemins
(Chapitre~\ref{chap:injections}) \\
Logique métier & Enchaînements non prévus, contournement d'étapes \\
\end{tableaucours}
\vspace{-0.8em}
\begin{center}{\small\itshape Tableau~14.1. Catégories du WSTG retenues pour LiZbiZ}\end{center}

% ============================================================================
\section{Cartographier}
\label{sec:c14-carto}
% ============================================================================

\begin{boxprocedure}
\textbf{Cartographier une application web}
\begin{enumerate}
    \item Parcourir l'application avec chaque rôle, en enregistrant les
          requêtes dans un relais local autorisé.
    \item Lister les pages, les points d'entrée et leurs paramètres.
    \item Noter les en-têtes, les cookies, les technologies détectées.
    \item Repérer les paramètres qui désignent un objet (numéro de colis,
          identifiant de compte).
    \item Dresser la matrice rôle par page, pour préparer les tests de
          contrôle d'accès.
\end{enumerate}
\end{boxprocedure}

\begin{boxalert}
\textbf{Violation typique.} Lancer un robot d'exploration agressif sur la
pré-production sans prévenir.

\smallskip
\textbf{Règle~:} l'exploration automatique crée des données, envoie des
courriels, déclenche des traitements. On fixe la cadence avec le client, on
évite les formulaires d'envoi, on prévient l'équipe d'exploitation.
\end{boxalert}

% ============================================================================
\section{Authentification et Session}
\label{sec:c14-auth}
% ============================================================================

\begin{tableaucours}{L{4.4cm}L{7.8cm}}{%
    \tableauHeader{Point testé} &
    \tableauHeader{Faiblesse recherchée}
}
Robustesse de la connexion & Verrouillage absent, énumération des comptes
par les messages \\
Réinitialisation de mot de passe & Jeton prévisible, lien qui n'expire
pas \\
Cookie de session & Absence des attributs de sécurité (indisponible au
script, réservé au canal chiffré) \\
Expiration & Session valide trop longtemps, pas d'invalidation à la
déconnexion \\
Second facteur & Absent sur les comptes sensibles \\
\end{tableaucours}
\vspace{-0.8em}
\begin{center}{\small\itshape Tableau~14.2. Tests d'authentification et de session}\end{center}

\begin{boxscenario}{Opération LOGONE. Le portail qui ne ferme pas la session}
\textbf{Constatations.} Sur le portail de suivi, le cookie de session n'a
pas l'attribut qui le réserve au canal chiffré~; à la déconnexion, le jeton
reste valide 8~heures. Un message de connexion distingue « compte inconnu »
de « mot de passe erroné », ce qui permet d'énumérer les comptes.

\medskip
\textbf{Analyse} (confiance élevée). Aucune de ces faiblesses n'est critique
seule~; ensemble, elles facilitent le vol et la réutilisation de session.
Elles alimentent le constat T-14 (délai de session trop long) du rapport.
\end{boxscenario}

% ============================================================================
\section{Lire un Historique HTTP}
\label{sec:c14-historique}
% ============================================================================

L'essentiel du test web se lit dans l'historique des requêtes et des
réponses. Le lab HE2-LAB-08 fournit un historique exporté~; l'élève y
cherche les points d'entrée, les paramètres d'objet, les en-têtes
manquants et les codes de réponse révélateurs.

\begin{tableaucours}{L{3.0cm}L{9.2cm}}{%
    \tableauHeader{Indice} &
    \tableauHeader{Ce qu'il peut révéler}
}
Paramètre numérique d'objet & Référence directe, à tester pour l'accès
indirect (Chapitre~\ref{chap:injections}) \\
Code 500 & Erreur serveur, message parfois bavard \\
En-tête de version & Technologie et version, utile à l'attaquant \\
Absence d'en-têtes de sécurité & Protection du navigateur non activée \\
Réponse différente selon le compte & Énumération possible \\
\end{tableaucours}
\vspace{-0.8em}
\begin{center}{\small\itshape Tableau~14.3. Lire un historique HTTP}\end{center}

% ============================================================================
\section{Classer selon le Top 10}
\label{sec:c14-top10}
% ============================================================================

Le \termecle{Top~10} de l'OWASP classe les grandes familles de risques des
applications web~\cite{owasptop10}. Il sert à ranger les constats et à
parler un langage commun avec les équipes de développement.

\begin{boxobjet}{T-12 · En-têtes de sécurité HTTP absents}
\textbf{Description~:} le portail de suivi ne renvoie pas les en-têtes de
sécurité recommandés (politique de contenu, type de contenu strict,
transport strict). \textbf{Preuve~:} réponses relevées dans l'historique.
\textbf{Gravité~:} faible. \textbf{Impact~:} protections du navigateur non
activées~; aggrave d'autres failles. \textbf{Recommandation~:} ajouter les
en-têtes au niveau du serveur ou de l'application. \textbf{Détection~:}
contrôle périodique des en-têtes de réponse.
\end{boxobjet}

\begin{boxmetier}{Pentesteur et développeur}
\textbf{Rapprocher le constat du code.} Le pentesteur décrit la requête qui
démontre la faille~; le développeur retrouve, à partir d'elle, la ligne de
code responsable. Un bon constat web contient toujours la requête exacte.
\end{boxmetier}

\separateur

\begin{boxresume}
\textbf{Ce qu'il faut retenir de ce chapitre~:}
\begin{itemize}
    \item \textbf{Tout ce qui vient du client} est contrôlable par le
          client.
    \item \textbf{Le WSTG} structure la campagne en catégories.
    \item \textbf{Cartographier d'abord}~: pages, paramètres, rôles.
    \item \textbf{Authentification et session}~: verrouillage, cookies,
          expiration.
    \item \textbf{L'historique HTTP} contient l'essentiel des indices.
    \item \textbf{Le Top~10} range les constats et parle aux développeurs.
\end{itemize}
\end{boxresume}

\begin{boxinfo}
\textbf{Points de contrôle de l'élève.} Je sais~: décrire requête, réponse
et session~; choisir des catégories du WSTG~; cartographier une
application~; tester authentification et session~; lire un historique
HTTP~; classer un constat dans le Top~10.
\end{boxinfo}

\bigskip

\begin{boxexo}[title={\ding{45}\quad Exercice~14.1. Confiance au client \niveaubase}]
Pourquoi un contrôle fait uniquement en JavaScript ne protège-t-il pas le
serveur~?
\end{boxexo}

\begin{boxexo}[title={\ding{45}\quad Exercice~14.2. Catégories \niveaubase}]
Citez trois catégories du WSTG et un test pour chacune.
\end{boxexo}

\begin{boxexo}[title={\ding{45}\quad Exercice~14.3. Cookie \niveaubase}]
À quoi servent les attributs « réservé au canal chiffré » et « indisponible
au script » d'un cookie de session~?
\end{boxexo}

\begin{boxexo}[title={\ding{45}\quad Exercice~14.4. Énumération \niveaumoyen}]
Pourquoi un message de connexion qui distingue « compte inconnu » de « mot
de passe erroné » est-il une faiblesse~? Quel message proposer~?
\end{boxexo}

\begin{boxexo}[title={\ding{45}\quad Exercice~14.5. Cartographie \niveaumoyen}]
Décrivez la matrice rôle par page du portail de suivi, avec les deux comptes
de test.
\end{boxexo}

\begin{boxexo}[title={\ding{45}\quad Exercice~14.6. Exploration prudente \niveaumoyen}]
Quelles précautions prenez-vous avant de lancer une exploration automatique
sur la pré-production~?
\end{boxexo}

\begin{boxexo}[title={\ding{45}\quad Exercice~14.7. Historique \niveaumoyen}]
Dans un historique, quels éléments trahissent un paramètre d'objet à
tester~?
\end{boxexo}

\begin{boxexo}[title={\ding{45}\quad Exercice~14.8. Session \niveaumoyen}]
Rédigez le constat du cookie de session sans attribut de canal chiffré et
sans invalidation à la déconnexion.
\end{boxexo}

\begin{boxexo}[title={\ding{45}\quad Exercice~14.9. Plan WSTG \niveauexpert}]
Établissez le plan de test, en catégories du WSTG, pour les 48~heures
accordées à M.~TAMO.
\end{boxexo}

\begin{boxexo}[title={\ding{45}\quad Exercice~14.10. Top 10 \niveauexpert}]
Rattachez chacun des constats web de l'Opération LOGONE à une catégorie du
Top~10, en justifiant.
\end{boxexo}

\bigskip

\begin{boxinfo}
\textbf{Pour aller plus loin.}
\begin{itemize}
    \item OWASP WSTG~\cite{owaspwstg}~; OWASP Top~10~\cite{owasptop10}.
\end{itemize}
\end{boxinfo}

\begin{boxinfo}
\textbf{Transition.} La catégorie la plus riche du test web est celle des
entrées mal validées~: le Chapitre~\ref{chap:injections} traite les
injections et les défauts de contrôle d'accès.
\end{boxinfo}
